% !TeX root = ../main.tex
\chapter{Planteamiento del problema}\label{cap:planteamiento}

Define con precisión el problema, los requisitos y la solución propuesta
a alto nivel, antes de entrar en detalles de implementación.

\section{Requisitos}

\begin{table}[H]
  \centering
  \caption{Requisitos del sistema.}
  \label{tab:requisitos}
  \begin{tabularx}{\linewidth}{llX}
    \toprule
    ID & Tipo & Descripción \\
    \midrule
    R1 & Funcional    & El sistema debe \ldots \\
    R2 & No funcional & El tiempo de respuesta debe ser inferior a \qty{10}{\milli\second}. \\
    \bottomrule
  \end{tabularx}
\end{table}

\section{Arquitectura propuesta}

Los diagramas de bloques sencillos se pueden dibujar directamente con
TikZ (\cref{fig:arquitectura}), de modo que el texto de la figura use la
misma tipografía que la memoria.

\begin{figure}[H]
  \centering
  \begin{tikzpicture}[
      node distance=1.8cm,
      bloque/.style={draw, rounded corners, minimum width=2.6cm,
                     minimum height=1cm, fill=blue!8, align=center},
      flecha/.style={-Stealth, thick}]
    \node[bloque] (entrada) {Adquisición};
    \node[bloque, right=of entrada] (proceso) {Procesado};
    \node[bloque, right=of proceso] (salida) {Visualización};
    \draw[flecha] (entrada) -- node[above]{\small datos} (proceso);
    \draw[flecha] (proceso) -- node[above]{\small resultados} (salida);
  \end{tikzpicture}
  \caption{Arquitectura general de la solución.}
  \label{fig:arquitectura}
\end{figure}

\section{Modelo matemático}

Las ecuaciones se numeran y referencian igual que las figuras
(\cref{ec:modelo}):
\begin{equation}
  \label{ec:modelo}
  y(t) = y_0 \, e^{-t/\tau} + \int_0^t h(t-s)\, u(s) \dif s
\end{equation}
donde $\tau = \qty{2.5}{\second}$ es la constante de tiempo.
